\chapter{Excepciones aceptadas y lo que falta medir}\label{cap:banco}

Ningún diseño real cumple el cien por ciento de las reglas que uno se pone. Lo importante es que cada excepción sea
\emph{consciente}: que se sepa qué regla se rompe, cuánto, por qué no había otra salida y quién la aceptó. Este
capítulo junta en un solo lugar todas las excepciones de la placa y, después, todo lo que no se puede verificar en el
computador y queda para el banco de laboratorio.

En este capítulo vas a ver:
\begin{itemize}
  \item la diferencia entre una excepción aceptada y una violación escondida;
  \item la lista completa de excepciones, con la regla, el valor y el motivo;
  \item las desviaciones del plan que no rompen ninguna regla;
  \item las comprobaciones en la vista previa de JLCPCB antes de pagar;
  \item la lista de mediciones de banco, con el valor esperado y cómo medir cada una.
\end{itemize}

\begin{idea}[excepción no es lo mismo que exclusión]
KiCad permite marcar una violación del DRC como ``excluida'' para que deje de aparecer. En esta placa \textbf{no se
excluyó ninguna violación ni se bajó ninguna severidad}: el DRC da 0 con todas las reglas activas
(capítulo~\ref{cap:verif}). Las excepciones de este capítulo son de otro tipo: casos en que una regla \emph{propia}
(del chequeo \archivo{revisar_pistas.py} o del plan de ruteo) no se cumple, el dueño lo revisó y lo aceptó. Cada una
queda escrita en el código del chequeo (tablas \texttt{ACEPTADOS}, \texttt{ACEPTADOS\_C} y
\texttt{DESACOPLO\_LIMITE}) con su motivo, y el chequeo las lista aparte en cada corrida.
\end{idea}

\section{Excepciones aceptadas}\label{sec:c12-excepciones}

\begin{longtable}{@{}L{0.22\linewidth}L{0.18\linewidth}L{0.17\linewidth}L{0.35\linewidth}@{}}
\caption{Lista completa de excepciones aceptadas por el dueño.}\label{tab:c12-excepciones}\\
\toprule
\textbf{Qué} & \textbf{Regla} & \textbf{Queda así} \medido & \textbf{Motivo y cuándo}\\\midrule\endfirsthead
\toprule
\textbf{Qué} & \textbf{Regla} & \textbf{Queda así} & \textbf{Motivo y cuándo}\\\midrule\endhead
C36, C46 (VCC del LMR36015, \SI{2.2}{\micro\farad}) & desacoplo a $\le\SI{1.0}{\milli\meter}$ del pin & \SI{3.00}{\milli\meter} &
La pista de VCC sale del pin con \SI{0.3}{\milli\meter}, la vía de tierra está pegada a C36.2 y el retorno al pin AGND va por \capa{F.Cu}. No hay lugar más cerca sin romper el lazo de potencia. Bloque 1.\\
C2, C3, C49, C50 (entrada de los LT3092) & $\le\SI{1.0}{\milli\meter}$ & \SI{4.30}{}, \SI{3.32}{}, \SI{4.04}{} y \SI{4.33}{\milli\meter} &
Se dejaron en la ubicación del equipo. La entrada es el riel de \SI{24}{\volt} protegido, que es DC. Bloque 3.\\
C12 (\SI{100}{\nano\farad}, \net{3V3} del PCM1808) & $\le\SI{1.0}{\milli\meter}$ & \SI{2.12}{\milli\meter} &
Se intercambió con C11 para que el \SI{100}{\nano\farad} quede más cerca; bajo el pin 4 no hay más lugar. Bloque 4.\\
C16 (\SI{100}{\nano\farad}, VREF del PCM1808) & $\le\SI{1.0}{\milli\meter}$ & \SI{1.32}{\milli\meter} &
Se intercambió con C15 por el mismo motivo. Bloque 4.\\
C11, C15, C18 (\SI{10}{\micro\farad} del PCM1808) & $\le\SI{1.0}{\milli\meter}$ & \SI{3.89}{}, \SI{1.52}{} y \SI{5.42}{\milli\meter} &
El \SI{100}{\nano\farad} de cada pin va primero (C12, C16 y C17); el de \SI{10}{\micro\farad} trabaja a frecuencias más bajas y tolera más distancia. Bloque 4, confirmado tras el bloque 5.\\
C5, C6, C52, C53 (desacoplo de los MCP4725) & se juzgan con $\le\SI{4}{\milli\meter}$ en vez de \SI{1}{\milli\meter} & \SIrange{1.22}{1.48}{\milli\meter} &
Es el límite del propio datasheet \doc{DS22039D §8.3}. Bloque 5.\\
Enlace C56.1 $\to$ R18.1 y C57.1 $\to$ R41.1 (nodo \net{SUM}) & entrar al pad por su eje, con el quiebre fuera del pad (chequeo C) & quiebre a \SI{45}{\degree} dentro del cobre de los pads; enlace de \SI{2.2}{\milli\meter} &
Los pads están desplazados \SI{1.52}{\milli\meter} en x y \SI{1.76}{\milli\meter} en y, con el otro pad del CF al lado. No hay forma geométrica de entrar por el eje; este es el enlace más corto, que es lo que importa en un nodo de alta impedancia. Bloque 5.\\
Pista de \SI{24}{\volt} de CH2 bajo J2 & nada bajo el cuerpo de un conector & cruza en y = \num{153.5}, a \SI{1.1}{\milli\meter} de los pines &
Es el único camino en \capa{F.Cu}. El bloque del BNC es de plástico PBT y queda a \SI{0.5}{\milli\meter} de la placa; el metal no la toca (plano rev.~B 2020 y fotos de LCSC). Bloque 3.\\
\net{IEPE_SV} bajo J8 & ídem & y = \num{153.825}; \SI{1.45}{\milli\meter} a J8.2 &
El mismo caso que J2. Bloque 6.\\
C19 (\SI{10}{\nano\farad}, filtro de VINR) & vía de tierra propia junto al pad (chequeo K) & sin vía; lo une el relleno de tierra de \capa{F.Cu} &
No hay lugar para una vía taponada a $\ge\SI{0.35}{\milli\meter}$ de los pads. Bloque 4.\\
L1, L2, COUT, R33, C47 & cuerpo a $\ge\SI{2.5}{\milli\meter}$ del borde sin rieles (JLC) & \SIrange{1.30}{1.85}{\milli\meter} &
Los rieles de montaje los agrega JLC según sus términos. Punto de control 2.\\
ERC & --- & 4 chequeos ignorados de fábrica &
Activados en una copia solo muestran 3 \texttt{four\_way\_junction} correctas (D1.K, \#PWR006, \#PWR018). Se dejaron como están.\\
\bottomrule
\end{longtable}

\begin{idea}[por qué el desacoplo tolera algo de distancia en esta placa]
Un condensador de desacoplo entrega la corriente rápida que pide el integrado. Cada milímetro de pista entre el
condensador y el pin agrega inductancia, y esa inductancia limita qué tan
rápido puede entregar corriente. Los integrados de esta placa son de audio y de baja velocidad (op-amps de
\SI{4.5}{\mega\hertz}, un ADC de \SI{48}{\kilo\hertz}, DACs lentos), así que unos milímetros pesan mucho menos que en
un procesador rápido. Por eso las excepciones se aceptaron en los condensadores ``grandes'' y en entradas de DC,
mientras que todos los desacoplos que sí importan (los CHF de los bucks, C7, C54, C17, C13) quedaron a menos de
\SI{1}{\milli\meter}.
\end{idea}

\section{Desviaciones del plan que no rompen ninguna regla}\label{sec:c12-desviaciones}

Estas no son excepciones a una regla, pero sí cambios respecto de lo que se había planeado. Se informaron al dueño en
cada bloque:

\begin{itemize}
  \item \textbf{Pistas de potencia más angostas que el valor por defecto de la clase}: \SIrange{0.25}{0.6}{\milli\meter}
  en lugar de \SI{0.5}{\milli\meter}, y vías de \num{0.6}/\SI{0.3}{\milli\meter} en lugar de \num{0.8}/\SI{0.4}{\milli\meter}.
  Las corrientes son de \SIrange{5}{20}{\milli\ampere}, con \SI{0.28}{\ampere} de pico en SW. Cada ancho se justifica
  con IPC-2221 en el capítulo~\ref{cap:calculos}.
  \item \textbf{Alimentación en \capa{In2.Cu} a \SI{0.4}{\milli\meter}} en vez de los \SI{0.5}{\milli\meter} previstos.
  Las caídas calculadas con los largos medidos son \SI{4.47}{\milli\volt} en \net{3V3} y \SI{1.75}{\milli\volt} en
  \net{5V_A}.
  \item \textbf{Tres puentes cortos en \capa{B.Cu} bajo el bus de \SI{24}{\volt}}: el tronco de \net{3V3}
  (\SI{2.2}{\milli\meter}) y los sensados de FB de U6 y U7 (\SI{2.15}{\milli\meter}). Se aprobaron en el bloque 1,
  cuando la placa era de 2 capas. Con \capa{In1.Cu} como plano entero ya no cortan ningún retorno.
  \item \textbf{Política de capas con 4 capas}: las redes críticas solo en \capa{F.Cu}, salvo \net{DC_MON*} e
  \net{IEPE_FEED*}. I2C de CH1, \net{DC_MON} y \net{DC_MON2} van en \capa{B.Cu}; I2C de CH2 y un tramo de
  \net{DC_MON2} en \capa{In2.Cu}. \archivo{cruces_capas.py} da 0 cruces bajo los nodos protegidos.
  \item \textbf{Cambios al chequeo propio}: \net{IEPE_FEED} se trata como alimentación (lo pidió el dueño); el anillo
  libre de un pad THT o vía propia cuenta como información en el chequeo I; el J pasó a mirar todas las capas; se
  agregó el M (serigrafía sobre vías).
  \item \textbf{Serigrafía}: a \SI{1.0}{\milli\meter} de alto no caben 21 designadores; quedan en
  \path{fabricacion/montaje.pdf}. Nunca se omiten marcas de polaridad o de pin 1, conectores ni puntos de prueba.
  \item \textbf{Ajustes por la serigrafía}: el punto de pin 1 del footprint RNX (U6/U7) pasó a la posición
  $(\num{-1.75};\ \num{-1.45})\,\si{\milli\meter}$ respecto del centro de la pieza. Las vías de tierra laterales de D2
  bajaron \SI{0.55}{\milli\meter} y las vías de \net{3V3} de U4 se corrieron.
\end{itemize}

\section{Antes de pagar: la vista previa de JLCPCB}\label{sec:c12-dfm}

Cuando subes los archivos, JLCPCB muestra una vista previa con las piezas colocadas sobre la placa (su revisión DFM).
Es la última oportunidad de atrapar un error de rotación, porque la colocación real la decide la biblioteca de montaje
de JLC, no nuestro cálculo (capítulo~\ref{cap:verif}). Hay que mirar:

\begin{enumerate}
  \item el pin 1 de U1 a U10 sobre la marca de pin 1 de la placa;
  \item el cátodo de D1 y D2 del lado de la franja de la serigrafía;
  \item el pin 1 del oscilador Y1;
  \item el centrado de U6 y U7 (encapsulado RNX): su footprint de JLC está girado \SI{90}{\degree} y corrido
  \SI{0.14}{\milli\meter} respecto del de KiCad, así que es el caso más delicado.
\end{enumerate}

\begin{ojo}[dos cosas del pedido que siguen abiertas]
\begin{itemize}
  \item \textbf{Stock}: el 2026-11-16, L2 (XAL4040-103MEC, C5144883) tenía solo 4 unidades. Hay que repetir
  \archivo{tools/stock_bom.py} el día del pedido.
  \item \textbf{Cantidad de placas ensambladas con Economic}: la idea es ensamblar 4 de 5. El cotizador pide los
  Gerber antes de mostrar esa opción, y en su demo Economic aparece deshabilitado (con Standard deja elegir de 2 a 5).
  Si Economic solo ofrece 2 o 5, se ensamblan 2. \archivo{fabricacion/LEEME.md} todavía dice ``hasta 4 placas''.
\end{itemize}
También hay que confirmar en el formulario la serigrafía blanca y las vías taponadas, que la página de capacidades no
menciona para Economic. Lo demás (4 capas, \SI{1.6}{\milli\meter}, verde, HASL sin plomo) sí aparece en su tabla.
\end{ojo}

\section{Lo que falta medir en el banco}\label{sec:c12-banco}

Hay cosas que ningún programa puede confirmar: el consumo real de un buck sin carga, la polarización de un micrófono
del que no hay datasheet o si un op-amp oscila con el cableado real. La lista que sigue está en el orden recomendado:
primero lo que puede dañar algo, después lo que depende de que lo anterior funcione.

\begin{ojo}[antes de energizar]
\begin{itemize}
  \item Con el multímetro en resistencia, comprueba que no haya corto entre \net{24V_PROT}, \net{3V3}, \net{5V} y
  \net{GND} (los pads de los condensadores de cada riel sirven de punto de medida).
  \item Fuente de banco a \SI{24}{\volt} con el límite de corriente en \SIrange{30}{50}{\milli\ampere}. Nunca más de
  \SI{26}{\volt} continuos: es la tensión de trabajo del TVS D2 \copia.
  \item Orden de encendido: primero el USB de la NUCLEO, después los \SI{24}{\volt} (cambio C12).
\end{itemize}
\end{ojo}

\begin{longtable}{@{}L{0.20\linewidth}L{0.30\linewidth}L{0.42\linewidth}@{}}
\caption{Mediciones de banco pendientes. Los puntos de prueba son: TP1 = \net{IEPE}, TP2 = \net{REST_OUT},
TP3 = \net{IEPE_SV}, TP4 = \net{REST_SV}, TP5--TP7 = \net{GND}.}\label{tab:c12-banco}\\
\toprule
\textbf{Qué} & \textbf{Valor esperado} & \textbf{Cómo medirlo}\\\midrule\endfirsthead
\toprule
\textbf{Qué} & \textbf{Valor esperado} & \textbf{Cómo medirlo}\\\midrule\endhead
Consumo de los bucks sin carga (FPWM) &
$\le\SI{0.15}{\watt}$ por integrado, unos \SI{6}{\milli\ampere} a \SI{24}{\volt} cada uno; $\le\SI{11}{\degreeCelsius}$ de subida \estimado &
Amperímetro de la fuente de banco con J1 y J7 abiertos (fuentes IEPE desconectadas). La lectura incluye las cargas de \net{3V3} y \net{5V}, así que es una cota del consumo de los dos bucks juntos. Es el dato menos conocido de la placa (capítulo~\ref{cap:calculos}).\\
Rieles de salida &
\net{3V3}: \SI{3.3}{\volt} $\pm\SI{1.5}{\percent}$; \net{5V}: \SI{5.0}{\volt} $\pm\SI{1.5}{\percent}$ \doc{SNVSB49D p.8}; \net{5V_A}: unos \SI{140}{\milli\volt} menos que \net{5V} con la carga máxima &
Multímetro en DC en los pads de un COUT (C31--C34 para \net{3V3}, C41--C44 para \net{5V}) y en C47 para \net{5V_A}, con la punta negra en TP5--TP7.\\
Rizado de los rieles &
componente a $\approx\SI{1}{\mega\hertz}$ (la frecuencia del buck), pequeña en \net{5V_A} por el filtro R33/C47 &
Osciloscopio en acople AC sobre los mismos pads, con la tierra de la punta lo más corta posible (resorte, no el cable largo con caimán).\\
Corriente de las fuentes IEPE &
CH1: \SI{2.81}{\milli\ampere}, y \SI{3.98}{\milli\ampere} con JP1 puesto. CH2: \SI{1.51}{\milli\ampere}, y \SI{2.51}{\milli\ampere} con JP2 &
Amperímetro en lugar del puente de J1 (CH1) o J7 (CH2): son borneras en serie pensadas para esto. Con un micrófono conectado o una resistencia de prueba de unos pocos k\si{\ohm} en el BNC; sin carga, la fuente se satura.\\
Polarización de cada micrófono &
RG-50: \SI{12.5}{\volt} \doc{ROGA RG-50}; SV 12: desconocida (normal entre \SIrange{7}{14}{\volt}). Cerca de \SI{0}{\volt} = corto; cerca de la alimentación = abierto &
Multímetro en DC en TP1 (CH1) y TP3 (CH2) con el micrófono conectado.\\
Temperatura del LT3092 en corto &
\SI{93}{\milli\watt}, unos $+\SI{5.3}{\degreeCelsius}$ ($\theta_{JA}$ \copia) &
Cortocircuitar el BNC un momento y medir la temperatura de U1/U8 (termómetro infrarrojo o termocupla). Confirma que un corto en el cable no daña nada.\\
Lazo de DC y TLV9152 &
\net{REST_OUT} en el objetivo que fije el firmware; se recomienda \SI{1.5}{\volt} (pendiente de aprobación). \net{VSHIFT} $= (\text{REST} + k\,V_B)/(1+k)$, con $k = \num{0.0997}$: debe quedar por debajo de unos \SI{2.7}{\volt} &
Multímetro en TP2 y TP4; \net{VSHIFT} se calcula con REST y la polarización medida en TP1/TP3 (o se mide con cuidado en el pin 7 de U3/U9).\\
Estabilidad del sumador y del seguidor &
sin oscilación, sin sobrepico al dar un escalón &
Osciloscopio en TP2/TP4. Escalón de corriente con JP1/JP2 y escalón del DAC por firmware. Mira también con la base de tiempo rápida (\si{\micro\second}/div) por si hay una oscilación de alta frecuencia.\\
Oscilador Y1 &
\SI{24.576}{\mega\hertz}; consumo $\le\SI{10}{\milli\ampere}$ \estimado &
Frecuencia con el osciloscopio en la salida (lado de R21). El consumo no se puede aislar: queda dentro del consumo total de \net{3V3}.\\
Identificación de los micrófonos &
SV 12 con conector TNC (IEPE), no LEMO de 7 pines; cápsula montada. ROGA: grabado RG-50, no MI17/MI19 &
Inspección visual antes de conectar. Si el SV 12 tiene LEMO de 7 pines, \textbf{no es IEPE}: no se conecta.\\
\bottomrule
\end{longtable}

\begin{ojo}[datos que no se miden pero tampoco están verificados]
Algunos datos de los cálculos vienen de copias y no del documento original: el cobre del barril de las vías
(\SI{18}{\micro\meter} \estimado), las fórmulas de IPC-2221 e IPC-2141, la $\theta_{JA}$ del LT3092 y los valores del
SMBJ26A ($V_\text{RWM}$ = \SI{26}{\volt}, $V_\text{BR}$ mín = \SI{28.9}{\volt}) \copia. Ninguno está cerca de un límite,
pero conviene contrastarlos con el datasheet cuando se tenga a mano.
\end{ojo}

\section*{Resumen}
\begin{itemize}
  \item No hay violaciones del DRC excluidas: las excepciones son reglas propias que el dueño aceptó, escritas en el
  código del chequeo con su motivo.
  \item Las excepciones son casi todas de distancia de desacoplo en piezas lentas o en entradas de DC, dos rutas bajo
  el plástico de los BNC y el enlace geométricamente imposible del nodo \net{SUM}.
  \item Antes de pagar: mirar en la vista previa el pin 1, los diodos, Y1 y los bucks, repetir el stock y confirmar la
  cantidad de placas ensambladas con Economic.
  \item En el banco, primero la seguridad (cortos, límite de corriente, orden de encendido), después rieles, corrientes
  IEPE, polarización de los micrófonos, lazo de DC y estabilidad.
\end{itemize}
